\subsection{Funtores hipertor locales de complejos de módulos cuasicoherentes: caso de los esquemas afines}
\label{subsection:III.6.4}
\label{III.6.4}

\begin{env}[6.4.1]
\label{III.6.4.1}
Sean $S$ un esquema afín de anillo $A$, y $X$ e $Y$ dos
$S$-esquemas afines de anillos $B$ y $C$, respectivamente; así,
$B$ y $C$ son $A$-álgebras. Todo complejo $\sh P_\bullet$ de
$\sh O_X$-módulos cuasicoherentes es de la forma
$\widetilde{P_\bullet}$ para un complejo $P_\bullet$ de
$B$-módulos, y análogamente todo complejo $\sh Q_\bullet$ de
$\sh O_Y$-módulos cuasicoherentes es de la forma
$\widetilde{Q_\bullet}$ para un complejo $Q_\bullet$ de
$C$-módulos (\sref[I]{I.1.3.7} y \sref[I]{I.1.3.8}).

Considerados como complejos de $A$-módulos, $P_\bullet$ y
$Q_\bullet$ tienen módulos hipertor
$\mathbf{Tor}_n^A(P_\bullet,Q_\bullet)$. Por bifuntorialidad, las
$A$-álgebras $B$ y $C$ actúan sobre este $A$-módulo; estas acciones
lo convierten en un $(B,C)$-bimódulo o, equivalentemente, en un
módulo sobre
\[
  B\otimes_AC=A(X\times_SY).
\]
Hemos definido así un $\sh O_{X\times_SY}$-módulo cuasicoherente
\[
  \shTor_n^{\sh O_S}(\sh P_\bullet,\sh Q_\bullet)
  =
  \bigl(\mathbf{Tor}_n^A(P_\bullet,Q_\bullet)\bigr)\supertilde.
  \tag{6.4.1.1}\label{III.6.4.1.1}
\]
Se denomina \emph{hipertor local de índice $n$} de
$\sh P_\bullet$ y $\sh Q_\bullet$, y también se denota por
$\shTor_n^S(\sh P_\bullet,\sh Q_\bullet)$.
\end{env}

\begin{lemma}[6.4.2]
\label{III.6.4.2}
Con las notaciones de \sref{III.6.4.1}, supongamos que
$A=R^{-1}A'$, donde $A'$ es un anillo y $R$ una parte
multiplicativa de $A'$. Sea $S'=\Spec(A')$. Entonces $X$ e $Y$
pueden considerarse como $S'$-preesquemas y
$X\times_{S'}Y=X\times_SY$
(\sref[I]{I.1.6.2} y \sref[I]{I.3.2.4}); además,
\[
  \shTor_\bullet^{S'}(\sh P_\bullet,\sh Q_\bullet)
  =
  \shTor_\bullet^S(\sh P_\bullet,\sh Q_\bullet).
\]
\end{lemma}

\begin{proof}
Esto se deduce de \sref{III.6.4.1.1} y \sref{III.6.3.13}.
\end{proof}

\begin{env}[6.4.3]
\label{III.6.4.3}
Con las notaciones e hipótesis de \sref{III.6.4.1}, sean
$\sh F=\widetilde F$ un $\sh O_X$-módulo cuasicoherente y
$\sh G=\widetilde G$ un $\sh O_Y$-módulo cuasicoherente.
Considerando $\sh F$ y $\sh G$ como complejos concentrados en
grado cero, denotamos por $\shTor_n^{\sh O_S}(\sh F,\sh G)$, o
por $\shTor_n^S(\sh F,\sh G)$, su hipertor local de índice $n$.
Por \sref{III.6.3.5}(ii),
\[
  \shTor_n^{\sh O_S}(\sh F,\sh G)
  =
  \bigl(\Tor_n^A(F,G)\bigr)\supertilde.
  \tag{6.4.3.1}\label{III.6.4.3.1}
\]

Volvamos ahora a dos complejos arbitrarios de módulos
cuasicoherentes $\sh P_\bullet$ y $\sh Q_\bullet$. Las fórmulas
\sref{III.6.4.1.1} y \sref{III.6.4.3.1}, junto con
\sref{III.6.3.2}, muestran que
$\shTor_\bullet^S(\sh P_\bullet,\sh Q_\bullet)$ es el objeto al
que convergen dos sucesiones espectrales
$'\sh E(\sh P_\bullet,\sh Q_\bullet)$ y
$''\sh E(\sh P_\bullet,\sh Q_\bullet)$, cuyos términos $E_2$ son
\[
  '\sh E_{pq}^2
  =
  \sh H_p\bigl(\shTor_q^S(\sh P_\bullet,\sh Q_\bullet)\bigr)
  \tag{6.4.3.2}\label{III.6.4.3.2}
\]
y
\[
  ''\sh E_{pq}^2
  =
  \bigoplus_{q'+q''=q}
  \shTor_p^S
  \bigl(\sh H_{q'}(\sh P_\bullet),
        \sh H_{q''}(\sh Q_\bullet)\bigr).
  \tag{6.4.3.3}\label{III.6.4.3.3}
\]
Aquí $\shTor_q^S(\sh P_\bullet,\sh Q_\bullet)$ en
\sref{III.6.4.3.2} denota el bicomplejo de
$\sh O_{X\times_SY}$-módulos cuasicoherentes
$\shTor_q^S(\sh P_i,\sh Q_j)$.
\end{env}

\begin{env}[6.4.4]
\label{III.6.4.4}
Sean
\[
  X^{(1)}=\Spec(B^{(1)}),
  \qquad
  Y^{(1)}=\Spec(C^{(1)})
\]
otros dos esquemas afines, donde $B^{(1)}$ y $C^{(1)}$ son
$A$-álgebras, y sean
\oldpage[III]{147}
\[
  u:X^{(1)}\longrightarrow X,
  \qquad
  v:Y^{(1)}\longrightarrow Y
\]
$S$-morfismos correspondientes a los $A$-homomorfismos
$\vphi:B\longrightarrow B^{(1)}$ y
$\psi:C\longrightarrow C^{(1)}$. Consideremos los complejos
\[
  u^*(\sh P_\bullet)
  =
  (P_\bullet\otimes_BB^{(1)})\supertilde,
  \qquad
  v^*(\sh Q_\bullet)
  =
  (Q_\bullet\otimes_CC^{(1)})\supertilde.
\]
Los $A$-homomorfismos canónicos
\[
  P_\bullet\longrightarrow P_\bullet\otimes_BB^{(1)},
  \qquad
  Q_\bullet\longrightarrow Q_\bullet\otimes_CC^{(1)}
\]
dan, por funtorialidad, un $A$-homomorfismo
\[
  \mathbf{Tor}_\bullet^A(P_\bullet,Q_\bullet)
  \longrightarrow
  \mathbf{Tor}_\bullet^A
  (P_\bullet\otimes_BB^{(1)},
   Q_\bullet\otimes_CC^{(1)}).
\]
De nuevo por funtorialidad, se trata en realidad de un homomorfismo
de $(B\otimes_AC)$-módulos. Por tanto, define un
$(u\times_Sv)$-morfismo
\[
  \theta:
  \shTor_\bullet^S(\sh P_\bullet,\sh Q_\bullet)
  \longrightarrow
  \shTor_\bullet^S(u^*\sh P_\bullet,v^*\sh Q_\bullet),
  \tag{6.4.4.1}\label{III.6.4.4.1}
\]
y, en consecuencia, un homomorfismo de
$\sh O_{X^{(1)}\times_SY^{(1)}}$-módulos
\[
  \theta^\sharp:
  (u\times_Sv)^*
  \shTor_\bullet^S(\sh P_\bullet,\sh Q_\bullet)
  \longrightarrow
  \shTor_\bullet^S(u^*\sh P_\bullet,v^*\sh Q_\bullet).
  \tag{6.4.4.2}\label{III.6.4.4.2}
\]
Este es manifiestamente un morfismo de bi-$\partial$-funtores en
las categorías de módulos cuasicoherentes acotados inferiormente.
El homomorfismo \sref{III.6.4.4.2} no es necesariamente biyectivo;
sin embargo:
\end{env}

\begin{lemma}[6.4.5]
\label{III.6.4.5}
Con las notaciones de \sref{III.6.4.4}, si $u$ y $v$ son
inmersiones abiertas, entonces \sref{III.6.4.4.2} es biyectivo.
\end{lemma}

\begin{proof}
Identifiquemos $X^{(1)}$ e $Y^{(1)}$ con abiertos de $X$ e $Y$,
respectivamente. Son uniones de abiertos de la forma $\mathrm D(f)$
y $\mathrm D(g)$, donde $f\in B$ y $g\in C$; los preesquemas
inducidos sobre $\mathrm D(\vphi(f))$ y $\mathrm D(f)$, y asimismo
sobre $\mathrm D(\psi(g))$ y $\mathrm D(g)$, son isomorfos.

Basta demostrar el lema cuando $X^{(1)}=\mathrm D(f)$ e
$Y^{(1)}=\mathrm D(g)$. En efecto, una vez sabido esto, en el caso
general basta mostrar que la restricción de $\theta^\sharp$ a cada
$\mathrm D(f)\times_S\mathrm D(g)$ es un isomorfismo. Si
\[
  u_1:\mathrm D(f)\longrightarrow X^{(1)},
  \qquad
  v_1:\mathrm D(g)\longrightarrow Y^{(1)}
\]
son las inclusiones canónicas, esta restricción es
$(u_1\times_Sv_1)^*(\theta^\sharp)$. De
\sref{III.6.4.4} y \sref[0]{0.4.4.8} se deduce inmediatamente que,
al componerla con el homomorfismo canónico
\[
  (u_1\times_Sv_1)^*
  \shTor_\bullet^S(u^*\sh P_\bullet,v^*\sh Q_\bullet)
  \longrightarrow
  \shTor_\bullet^S((u')^*\sh P_\bullet,(v')^*\sh Q_\bullet),
  \tag{6.4.5.1}\label{III.6.4.5.1}
\]
donde $u'=u\circ u_1$ y $v'=v\circ v_1$, se obtiene el
homomorfismo canónico
\[
  (u'\times_Sv')^*
  \shTor_\bullet^S(\sh P_\bullet,\sh Q_\bullet)
  \longrightarrow
  \shTor_\bullet^S((u')^*\sh P_\bullet,(v')^*\sh Q_\bullet).
  \tag{6.4.5.2}\label{III.6.4.5.2}
\]
Si \sref{III.6.4.5.1} y \sref{III.6.4.5.2} son isomorfismos,
también lo es $(u_1\times_Sv_1)^*(\theta^\sharp)$.

Supongamos, pues, que $X^{(1)}=\mathrm D(f)$ e
$Y^{(1)}=\mathrm D(g)$, de modo que $B^{(1)}=B_f$ y
$C^{(1)}=C_g$. Entonces $u^*(\sh P_\bullet)$ y
$v^*(\sh Q_\bullet)$ se identifican, respectivamente, con
$((P_\bullet)_f)\supertilde$ y $((Q_\bullet)_g)\supertilde$.
Además, $X^{(1)}\times_SY^{(1)}$ se identifica con el subesquema
abierto $\mathrm D(f\otimes g)$ de
\[
  X\times_SY=\Spec(B\otimes_AC)
\]
(\sref[II]{II.4.3.2.4}). Queda por demostrar que el homomorfismo
\[
  \bigl(\mathbf{Tor}_\bullet^A(P_\bullet,Q_\bullet)\bigr)_{f\otimes g}
  \longrightarrow
  \mathbf{Tor}_\bullet^A((P_\bullet)_f,(Q_\bullet)_g),
  \tag{6.4.5.3}\label{III.6.4.5.3}
\]
\oldpage[III]{148}
inducido por los homomorfismos canónicos
$P_\bullet\longrightarrow(P_\bullet)_f$ y
$Q_\bullet\longrightarrow(Q_\bullet)_g$, es biyectivo.

Por \sref[0]{0.1.6.1},
\[
  (P_\bullet)_f=\varinjlim_m P_\bullet^{(m)},
\]
donde cada $P_\bullet^{(m)}$ es el mismo complejo de $B$-módulos
que $P_\bullet$, y para $m\leq n$ la aplicación
$P_\bullet^{(m)}\longrightarrow P_\bullet^{(n)}$ es la
multiplicación por $f^{\,n-m}$. El enunciado análogo vale para
$Q_\bullet$, sustituyendo $f$ por $g$. El homomorfismo
\[
  \mathbf{Tor}_\bullet^A(P_\bullet^{(m)},Q_\bullet^{(m)})
  \longrightarrow
  \mathbf{Tor}_\bullet^A(P_\bullet^{(n)},Q_\bullet^{(n)})
\]
correspondiente a estas aplicaciones de transición es, por
definición, la multiplicación por $(f\otimes g)^{n-m}$. La
conclusión se deduce de \sref[0]{0.1.6.1}, aplicado al miembro
izquierdo de \sref{III.6.4.5.3}, y de \sref{III.6.3.6}.
\end{proof}

\begin{env}[6.4.6]
\label{III.6.4.6}
Con las notaciones de \sref{III.6.4.4}, se definen del mismo modo
homomorfismos canónicos de funtores espectrales
\[
  \left\{
  \begin{aligned}
  (u\times_Sv)^*\,{}'\sh E(\sh P_\bullet,\sh Q_\bullet)
  &\longrightarrow
  {}'\sh E(u^*\sh P_\bullet,v^*\sh Q_\bullet),\\
  (u\times_Sv)^*\,{}''\sh E(\sh P_\bullet,\sh Q_\bullet)
  &\longrightarrow
  {}''\sh E(u^*\sh P_\bullet,v^*\sh Q_\bullet).
  \end{aligned}
  \right.
  \tag{6.4.6.1}\label{III.6.4.6.1}
\]
El razonamiento de \sref{III.6.4.5} muestra que, cuando $u$ y $v$
son inmersiones abiertas, estos homomorfismos son biyectivos.
Teniendo en cuenta \sref{III.6.3.6.2} y \sref{III.6.3.6.3}, el
razonamiento demuestra que son isomorfismos en los términos $E^2$,
mientras que \sref{III.6.4.5} demuestra que son isomorfismos en los
objetos límite. La afirmación se deduce entonces de
\sref[0]{0.11.1.2} y \sref[0]{0.11.2.4}.
\end{env}
